% ══════════════════════════════════════════════════════════════════════════════
%  EXPERIMENT 1 — WATER CYCLE IN A BAG                                  (Lesson 2)
% ══════════════════════════════════════════════════════════════════════════════
\experimentdiv{Water Cycle in a Bag}{Building a model \textbullet{} Observing over several days \textbullet{} Strengths and limitations of a model}{2}\label{sec:1}
\lessoncodes{VC2S8U09 \textbullet{} VC2S8I01, I03, I05, I06, I08}

\vocabhead
\begin{ctab}
\begin{tabular}{|SW{4.2cm}|W{11.6cm}|}
  \hline\hd & \thead{Word} & \thead{Definition (copy from the slide)} \\ \hline
  & \scaf{\textbf{Model}} & \scaf{A simpler version of something real, used to explain it or make predictions.} \\ \hline
  & \textbf{Evaporation} & \\ \hline
  & \textbf{Condensation} & \\ \hline
  & \textbf{Precipitation} & \\ \hline
  & \textbf{Limitation} & \\ \hline
\end{tabular}
\end{ctab}

\begin{lipanel}
\inquiryli{Make and record observations of a model over time, and evaluate the model.}
\understandli{Identify evaporation, condensation and precipitation in the model.}
\sctitle
\begin{itemize}[leftmargin=1.4em, itemsep=0pt, topsep=1pt]
  \item I can record clear observations of my model over several days.
  \item I can match each part of my model to a stage of the water cycle.
  \item I can give one strength and one limitation of the model.
\end{itemize}
\end{lipanel}

\begin{skillsfocus}
\focus{QP1}{Predict what will happen in the bag.}
\focus{PC6}{Observe the bag at each lesson.}
\focus{PA4}{Build a model of the water cycle.}
\focus{PA8}{Match the model to the real cycle.}
\focus{EV2}{Compare the sunny bag and the shaded bag.}
\focus{EV3}{Name the model's limitations.}
\focus{CO2}{Draw the set-up as a scientific diagram.}
\end{skillsfocus}

% ─── TECHNIQUE ────────────────────────────────────────────────────────────────
\exppart{Technique spotlight}{Blue}

\begin{technique}{Observing over time}{PC6}
Some experiments change slowly, so you observe them again and again. To make the observations comparable:
\begin{enumerate}[leftmargin=1.6em, itemsep=1pt, topsep=3pt]
  \item Observe at about the same time of day, from the same position.
  \item Record the date, time and weather each time.
  \item Write what you see in exact words: \emph{where}, \emph{how much}, \emph{what size}.
  \item Add a quick labelled sketch: a picture often shows what words miss.
\end{enumerate}
\end{technique}

\begin{workedex}{A precise observation}
\textbf{Vague:} \emph{The bag looked wet.}

\smallskip
\textbf{Precise:} \emph{Small clear droplets covered the top third of the bag. A few larger drops had run down the side facing the window.}
\end{workedex}

\begin{wedo}{A precise observation}
Rewrite this observation so that it is precise: \emph{There was some water at the bottom.}
\answerbox{3.5cm}
\end{wedo}

\begin{technique}{Strengths and limitations of a model}{PA8}
A \textbf{model} is useful because it shows something real in a simpler way. Every model also leaves things out or gets things wrong: these are its \textbf{limitations}. Judge a model by saying what it shows well (a strength) and what it does not show (a limitation).
\end{technique}

\begin{workedex}{A globe}
\textbf{Strength:} a globe shows that Earth is round, and where the continents and oceans are.

\smallskip
\textbf{Limitation:} it does not show weather or clouds, and its mountains are not to scale.
\end{workedex}

\begin{wedo}{A train-line map}
A map of Melbourne's train lines is a model. Give one strength and one limitation.
\answerbox{4.5cm}
\end{wedo}

\begin{youdobox}[Matching the model]
\task Each part of the bag model stands for something in the real water cycle. Complete the table. \skr{PA8}
\begin{ctab}
\begin{tabular}{|SW{5.0cm}|W{4.6cm}|W{5.4cm}|}
  \hline\hd & \thead{In the bag} & \thead{In the real world} & \thead{Stage of the water cycle} \\ \hline
  & \scaf{Water at the bottom} & \scaf{Oceans, lakes and rivers} & \scaf{Collection} \\ \hline
  & Sunlight through the window & & \\ \hline
  & Water vapour in the air of the bag & & \\ \hline
  & Droplets on the top of the bag & & \\ \hline
  & Drops running down the sides & & \\ \hline
\end{tabular}
\end{ctab}
\end{youdobox}

% ─── PLAN AND SAFETY / METHOD ────────────────────────────────────────────────
\exppart{Plan, safety and method}{Teal}

\textbf{Aim:} To model the water cycle in a sealed bag and observe evaporation, condensation and precipitation.

\needspace{10\baselineskip}
\task Predict what you will see inside a bag taped to a sunny window after one day, and inside a bag kept in a dark cupboard. Explain your prediction. \skr{QP1}
\answerbox{5cm}

\subsubsection*{Equipment}
Resealable plastic bag (about 20\,cm $\times$ 25\,cm) \textbullet{} 100\,mL measuring cylinder \textbullet{} water and blue food colouring \textbullet{} permanent marker \textbullet{} strong sticky tape \textbullet{} a sunny window (and a dark cupboard for the class's second bag)

\begin{center}
\begin{tikzpicture}[scale=0.78]
  \draw[eq] (0,0) -- (0,6);
  \draw[Ink, line width=0.4pt] (-0.25,0) -- (-0.25,6);
  \draw[eq] (0.15,0.6) rectangle (3.4,4.6);
  \draw[Ink, line width=1pt] (0.15,4.3) -- (3.4,4.3);
  \fill[black!12] (0.2,0.65) rectangle (3.35,1.4);
  \draw[Ink, line width=0.8pt] (0.2,1.4) -- (3.35,1.4);
  \draw[Ink, line width=1.6pt] (0.6,4.75) -- (1.4,4.45) (2.2,4.45) -- (3.0,4.75);
  \draw[eq] (2.6,3.4) circle (0.3);
  \draw[eq] (0.8,3.3) .. controls (0.8,3.7) and (1.6,3.7) .. (1.6,3.3) -- cycle;
  \foreach \p in {(0.6,3.95),(1.2,4.05),(1.9,3.95),(2.9,4.0),(3.15,2.6)}{\fill[Ink] \p circle (0.05);}
  \eqlab{(3.25,4.6)}{(4.6,5.2)}{zip seal}
  \eqlab{(2.9,4.0)}{(4.6,4.3)}{droplets}
  \eqlab{(2.85,3.45)}{(4.6,3.5)}{sun and cloud drawn in marker}
  \eqlab{(3.35,1.0)}{(4.6,1.0)}{100\,mL coloured water}
  \eqlab{(1.0,4.6)}{(1.4,5.6)}{sticky tape}
  \eqlabl{(-0.25,2.5)}{(-0.9,2.5)}{window}
\end{tikzpicture}
\expfigcap{Scientific diagram of the water cycle bag taped to a window.}
\end{center}

\begin{infobox}{Safety}
Seal the bag fully and tape it firmly so it cannot fall. Food colouring stains: wipe up drips straight away.
\end{infobox}

\subsubsection*{Method}
\begin{enumerate}[itemsep=2pt]
  \item With the marker, draw a sun and a cloud near the top of the bag, and waves near the bottom.
  \item Measure 100\,mL of water, add two drops of food colouring, and pour it into the bag.
  \item Squeeze out most of the air and seal the bag. Check it does not leak.
  \item Tape the bag to a sunny window, seal at the top. Your teacher puts a second bag in a dark cupboard.
  \item Observe both bags straight away and at the start of each Science lesson for a week. Record each observation in the table.
\end{enumerate}

% ─── RESULTS ─────────────────────────────────────────────────────────────────
\exppart{Results}{Amber}

\subsubsection*{Observations of the two bags \skr{PC6}}
\begin{ctab}
\begin{tabular}{|SW{2.6cm}|W{2.6cm}|W{5.0cm}|W{4.8cm}|}
  \hline\hd & \thead{Day and time} & \thead{Weather} & \thead{Bag on the sunny window} & \thead{Bag in the cupboard} \\ \hline
  & \scaf{e.g. Day 0, 9:10 am} & \scaf{Sunny, warm} & \scaf{No droplets; water clear blue} & \scaf{No droplets} \\ \hline
  & Day 0 & & & \\ \hline
  & Day 1 & & & \\ \hline
  & Day 2 & & & \\ \hline
  & Day 3 & & & \\ \hline
  & Day 4 & & & \\ \hline
\end{tabular}
\end{ctab}

\needspace{14\baselineskip}
\task Sketch the sunny bag on the first day and on the last day. Label the droplets, the water and any drops running down. \skr{CO2}

\noindent
\begin{minipage}[t]{0.48\linewidth}\drawbox{6cm}{}\nameline{Day 0}\end{minipage}\hfill
\begin{minipage}[t]{0.48\linewidth}\drawbox{6cm}{}\nameline{Last day}\end{minipage}

% ─── ANALYSIS ────────────────────────────────────────────────────────────────
\exppart{Analysis and evaluation}{Coral}

\task Describe where in the sunny bag you saw evidence of evaporation, condensation and precipitation. Quote your observations. \skr{PA8}
\answerbox{5cm}

\needspace{10\baselineskip}
\task Compare the sunny bag with the bag in the cupboard. Explain the difference, using the word \emph{energy}. \skr{EV2}
\answerbox{5cm}

\needspace{12\baselineskip}
\task Look at the droplets at the top of the sunny bag.
\pt{a}Were they blue, or clear?
\answerbox{1cm}
\pt{b}Explain why. Use this to explain why rain is fresh water even when it comes from the salty sea.
\answerbox{5cm}

\needspace{10\baselineskip}
\task The amount of water at the bottom of the bag hardly changed all week. Explain why, using the word \emph{cycle}.
\answerbox{5cm}

\needspace{14\baselineskip}
\task Evaluate the bag as a model of the water cycle. \skr{EV3}
\pt{a}Give one strength of the model.
\answerbox{5cm}
\pt{b}Give two limitations of the model.
\answerbox{5cm}

\begin{yr8box}
\begin{keepext}
\exttask Plants release water vapour from their leaves (transpiration). Design a change to the bag model that would show transpiration as well, and predict what you would observe.
\answerbox{5cm}
\end{keepext}
\begin{keepext}
\exttask Earth's water cycle has been recycling the same water for billions of years. Use the bag model to explain how this is possible.
\answerbox{5cm}
\end{keepext}
\end{yr8box}

% ─── SKILL CHECK ─────────────────────────────────────────────────────────────
\exppart{Skill check}{NavyBlue}

\begin{skillcheck}
\scrow{QP1}{I can predict and explain what will happen.}
\scrow{PC6}{I recorded precise observations at each lesson, with the date and weather.}
\scrow{PA4}{I built a working model of the water cycle.}
\scrow{PA8}{I can match each part of the model to a stage of the water cycle.}
\scrow{EV2}{I can compare the two bags and explain the difference.}
\scrow{EV3}{I can give a strength and a limitation of the model.}
\scrow{CO2}{I can draw and label the set-up as a scientific diagram.}
\end{skillcheck}

\task Which skill will you focus on in Experiment 2? Explain what you will do differently.
\answerbox{5cm}
